\chapter{Free QFT Review}

A free quantum field theory does not account for interactions between particles. It describes particles that do not interact with each other, and the fields that represent them evolve independently. Free QFTs are often used as a starting point for understanding more complex interacting theories, as they provide a simpler framework to study the properties of quantum fields and particles.

The lagrangian for a free quantum field theory typically consists of kinetic terms and mass terms for the fields, without any interaction terms: it can only contain \textbf{quadratic terms} in the fields, without any higher-order interactions.
Quanta of the free fields correspond to particles that do not interact with each other, and their dynamics are governed by the \textbf{free propagators} of the theory, which describe how particles propagate through spacetime without any interactions: propagators are the Green's functions of the free theory, and they determine how particles propagate from one point to another in spacetime with determined probability amplitudes.

We will review the free quantum field theories for different types of particles, including scalar fields (spin 0), fermionic fields (spin 1/2), and gauge fields (spin 1). We will also discuss the symmetries of these theories and how they lead to conservation laws through Noether's theorem. This review will provide a foundation for understanding more complex interacting quantum field theories, and it will also highlight the differences between free and interacting theories, as well as the role of symmetries in quantum field theory.

\section{Spin 0}

For a real scalar field $\phi(x)$, the free Lagrangian is given by
\[
    \mathcal{L} = \frac{1}{2} (\partial_\mu \phi)^2 - \frac{1}{2} m^2 \phi^2,
    \label{eq:free_KG_lagrangian}
\]
where $m$ is the mass of the scalar particle. The equation of motion derived from this Lagrangian is the \textbf{Klein-Gordon equation}:
\[
    (\Box + m^2) \phi = 0,
    \label{eq:free_KG_equation}
\]
where $\Box = \partial_\mu \partial^\mu$ is the d'Alembertian operator. The solutions to this equation describe free scalar particles with mass $m$. This is a classical equation of motion, having the same role of Maxwell's equations in classical electromagnetism, but for a scalar field \(\phi(x)\) instead of a vector field \(A^{\mu}(x)\).

The idea of quantum field theory is to promote the classical field \(\phi(x)\) to a quantum operator \(\hat{\phi}(x)\) that acts on a Hilbert space of states. The quantization procedure involves imposing commutation relations on the field operators, which leads to the creation and annihilation of particles. The free scalar field can be expanded in terms of creation and annihilation operators, which allows us to describe the particle content of the theory. Indeed, we can see how decomposing the field into its Fourier modes leads to the interpretation of the quanta of the field as particles with specific momenta and energies:
\[
    \phi(t,\, \mathbf{x}) = \int \frac{\mathrm{d}^3 \mathbf{k}}{(2\pi)^3} \, e^{i \mathbf{k} \cdot \mathbf{x}} \tilde{\phi}(t,\,\mathbf{k}),
\]
where each mode can be computed by using this ansatz in the Klein-Gordon equation in momentum space, leading to the dispersion relation of an harmonic oscillator represented by on-shell particles with energy \(E_{\mathbf{k}} = \omega_{\mathbf{k}} = \sqrt{\mathbf{k}^2 + m^2}\):
\[
    \left( \partial_t^2 + \omega_{\mathbf{k}}^2 \right) \tilde{\phi}(t, \mathbf{k}) = 0,
\]
showing how the free scalar field can be interpreted as a collection of independent harmonic oscillators, each corresponding to a different momentum mode.

In order to quantize the free scalar field, we have to promote the fields themselves to operators and impose the canonical commutation relations, which leads to the creation and annihilation of particles. Before showing how to do that, it is crucial to note that \(\hat{\phi}(x)\) depends on both space and time, enforcing the Heisenberg picture of quantum mechanics, where operators evolve in time while states remain fixed. The canonical commutation relations for the scalar field are given by
\[
    [\hat{\phi}(t, \mathbf{x}), \hat{\pi}(t, \mathbf{y})] = i \delta^3(\mathbf{x} - \mathbf{y}),
\]
where \(\hat{\pi}(x) = \frac{\partial \mathcal{L}}{\partial (\partial_t \hat{\phi}(x))} = \partial_t \hat{\phi}(x)\) is the conjugate momentum operator. This constraint along with reality of the field \(\hat{\phi}(x) = \hat{\phi}^\dagger(x)\) allows us to express the field operator in terms of \textbf{creation and annihilation operators}:
\[
    \hat{\phi}(x) = \int \frac{\mathrm{d}^3 \mathbf{k}}{(2\pi)^3} \frac{1}{\sqrt{2 \omega_{\mathbf{k}}}} \left( \hat{a}_{\mathbf{k}} e^{-i k \cdot x} + \hat{a}_{\mathbf{k}}^\dagger e^{i k \cdot x} \right),
\]
where \(k \cdot x = \omega_{\mathbf{k}} t - \mathbf{k} \cdot \mathbf{x} = k^{\mu} x_{\mu}\), and \(\omega_{\mathbf{k}} = k^0\). The creation and annihilation (ladder) operators satisfy the following commutation relations
\[
    \begin{dcases}
        [\hat{a}_{\mathbf{k}}, \hat{a}_{\mathbf{k}'}^\dagger] = (2\pi)^3 \delta^3(\mathbf{k} - \mathbf{k}'), \\
        [\hat{a}_{\mathbf{k}}, \hat{a}_{\mathbf{k}'}] = [\hat{a}_{\mathbf{k}}^\dagger, \hat{a}_{\mathbf{k}'}^\dagger] = 0.
    \end{dcases}
\]
Note that while the field operator \(\hat{\phi}(x)\) is Hermitian and time dependent, the creation and annihilation operators are not, and they are responsible for creating and destroying particles in the quantum field theory. The \textbf{vacuum state} \(\ket{0}\) is defined as the state that is annihilated by all annihilation operators:
\[
    \hat{a}_{\mathbf{k}} \ket{0} = 0 \quad \forall \mathbf{k},
\]
while we can create one-particle states by acting with the creation operators on the vacuum, which adds one quanta of the field with momentum \(\mathbf{k}\):
\[
    \ket{\mathbf{k}} = \sqrt{2 \omega_{\mathbf{k}}} \hat{a}_{\mathbf{k}}^\dagger \ket{0},
\]
and more generally, we can create multi-particle states by acting with multiple creation operators on the vacuum, which leads to the \textbf{Fock space} of the theory. With this definition of single particle states, we can compute the inner product between two single particle states, which is given by
\[
    \bra{\mathbf{k}} \ket{\mathbf{k}'} = 2 \omega_{\mathbf{k}} \delta^3(\mathbf{k} - \mathbf{k}'),
\]
showing that the single particle states are orthogonal and normalized with respect to the relativistic normalization convention. The normalization \(\sqrt{2 \omega_{\mathbf{k}}}\) is arbitrary, but important for ensuring that the probabilities computed in the theory are consistent with relativistic invariance.

\begin{remark}
    To find a Lorentz invariant measure, we can start from the Minkowski space measure \(\mathrm{d}^4 k\) and impose the on-shell condition \(k^2 = m^2\) using a delta function along with a step function to ensure that we only consider positive energy solutions:
    \[
        \int \mathrm{d}^4 k \, \delta(k^2 - m^2) \theta(k^0) = \int \frac{\mathrm{d}^3 \mathbf{k}}{2 \omega_{\mathbf{k}}}.
    \]
    We have then found the Lorentz invariant measure \(\frac{\mathrm{d}^3 \mathbf{k}}{2 \omega_{\mathbf{k}}}\) for integrating over the momentum space of on-shell particles, which is crucial for ensuring that the theory is consistent with special relativity. The previous normalization of the single particle states is consistent with this Lorentz invariant measure, since it ensures that the inner product between single particle states is Lorentz invariant, and it also allows us to compute probabilities and cross sections in a way that is consistent with relativistic invariance:
    \[
        \int \frac{\mathrm{d}^3 \mathbf{k}}{(2\pi)^3} \frac{1}{2 \omega_{\mathbf{k}}} \ket{\mathbf{k}} \bra{\mathbf{k}} = 1,
    \]
    which is the completeness relation for the single particle states in the theory. This ensures that we can express any state in the Fock space as a superposition of single particle states, and it also allows us to compute transition amplitudes and probabilities in a way that is consistent with relativistic invariance.
\end{remark}

\subsection{Feynman's Propagator}

In Quantum Field Theory (QFT), we can compute physical observables in terms of \textit{Green functions}, which are defined as the expectation values of products of field operators evaluated in the vacuum state:
\[
    \langle 0 | T\{ \hat{\phi}(x_1) \dots \hat{\phi}(x_n) \} | 0 \rangle
\]
where "$T$" stands for "time ordering" and defines the time-ordered product. For two bosonic field operators, the time-ordered product is explicitly given by:
\[
    T \{ \hat{O}(x) \hat{O}(y) \} = \Theta(x^0 - y^0) \hat{O}(x) \hat{O}(y) + \Theta(y^0 - x^0) \hat{O}(y) \hat{O}(x)
\]
where $\Theta(t)$ is the Heaviside step function, ensuring that operators evaluated at later times always appear to the left.

A very important case is the 2-point function, universally known as the \textbf{Feynman Propagator}. For a real scalar field, we define it as:
\[
    D_F(x, y) = D_F(x - y) \equiv \langle 0 | T \{ \hat{\phi}(x) \hat{\phi}(y) \} | 0 \rangle
\]
Notice that the propagator depends only on the difference $x-y$ due to space-time translation symmetry. Physically, it can be interpreted as the probability amplitude for a "causal" propagation of a particle between the spacetime points $x$ and $y$.

The exact analytical result for this expectation value (the proof can be found in standard textbooks such as Peskin 2.4 or Schwartz 6.2) is given by the following momentum-space integral:
\[
    D_F(x, y) = \int \frac{d^4k}{(2\pi)^4} \frac{i}{k^2 - m^2 + i\epsilon} e^{-ik \cdot (x - y)}
\]

\textbf{Comments:}
\begin{itemize}
    \item This result is strictly valid for a \textit{free theory} (non-interacting fields).
    \item Here, $k^0$ is an independent integration variable, which means the particle can be "off-shell", i.e., $k^0 \neq \sqrt{\vec{k}^2 + m^2}$.
    \item The "$i\epsilon$" term is known as the \textit{Feynman prescription}. Mathematically, it shifts the poles of the integrand at $k^2 = m^2$ slightly away from the real axis in the complex $k^0$-plane, avoiding singularities along the integration contour. This specific shift properly implements the causal time-ordering defined by the $\Theta$ functions. It is implied that $\epsilon \to 0^+$, and it is common to write the denominator simply as $k^2 - m^2 + i0^+$ or even leave it implicit.
\end{itemize}

Furthermore, $D_F(x, y)$ is mathematically a "Green function" of the classical Klein-Gordon equation. This means it acts as the inverse of the Klein-Gordon differential operator:
\[
    (\square_x + m^2) D_F(x, y) = -i \delta^{(4)}(x - y)
\]
The specific boundary conditions chosen to solve this differential equation are precisely what give rise to the aforementioned "$i\epsilon$" prescription in momentum space.


\subsection{Noether's Theorem}

Noether's Theorem establishes a profound connection between symmetries and conservation laws: a continuous symmetry of the action implies the existence of a \textit{conserved current} and a corresponding \textit{conserved charge}.

Consider a system described by a Lagrangian density $\mathcal{L} = \mathcal{L}(\phi_k, \partial_\mu \phi_k)$, where $\phi_k$ can be any type of field. Suppose there is an infinitesimal transformation of the fields $\phi_k$ and the coordinates $x^\mu$, described by a continuous parameter $\alpha$. If this transformation leaves the total action $S = \int d^4x \, \mathcal{L}(x)$ invariant, the variation of the Lagrangian density must at most be a total divergence:
\[
    \delta \mathcal{L} = \alpha \partial_\mu \mathcal{J}^\mu + \mathcal{O}(\alpha^2)
\]
for some 4-vector $\mathcal{J}^\mu$. Note that if $\mathcal{L}$ is strictly invariant point-by-point, then $\mathcal{J}^\mu = 0$.

By applying the Euler-Lagrange equations, we can extract the general form of the conserved Noether current $J^\mu$:
\[
    J^\mu = \sum_k \frac{\partial \mathcal{L}}{\partial(\partial_\mu \phi_k)} \frac{\partial \phi_k}{\partial \alpha} - \mathcal{J}^\mu
\]
which satisfies the continuity equation:
\[
    \partial_\mu J^\mu = 0
\]
This local conservation law implies the existence of a globally conserved charge $Q$, obtained by integrating the time-component of the current ($J^0$) over all space:
\[
    Q = \int d^3x \, J^0 \implies \frac{dQ}{dt} = 0
\]

\subsubsection{Energy-Momentum Tensor}

An exceptionally important symmetry is space-time translation: $x^\nu \to x^\nu + a^\nu$. Because the parameter $a^\nu$ has 4 independent components (for $\nu=0,1,2,3$), this symmetry yields 4 distinct conserved currents. These currents are naturally arranged into a rank-2 tensor known as the \textbf{Energy-Momentum Tensor}, $T^{\mu\nu}$, which satisfies $\partial_\mu T^{\mu\nu} = 0$.

Applying the general formula for $J^\mu$ to translations (where the variation of the field is $\delta \phi_k = a^\nu \partial_\nu \phi_k$ and $\mathcal{J}^\mu = g^{\mu\nu} \mathcal{L}$), we obtain:
\[
    T^{\mu\nu} = \sum_k \frac{\partial \mathcal{L}}{\partial(\partial_\mu \phi_k)} \partial^\nu \phi_k - g^{\mu\nu} \mathcal{L}
\]

A crucial component of this tensor is the \textit{energy density} of the field configuration, given by the $00$-component:
\[
    \mathcal{E} = T^{00}
\]
For Lagrangians where the canonical momentum is standardly defined as $\pi \equiv \frac{\partial \mathcal{L}}{\partial(\partial_t \phi)}$, this energy density perfectly coincides with the Hamiltonian density: $\mathcal{E} = \mathcal{H}$.

For a real scalar field, after performing canonical quantization (promoting fields to operators and imposing commutation relations), we obtain the full quantum Hamiltonian operator by integrating the Hamiltonian density over space:
\[
    \hat{H} = \int \frac{d^3k}{(2\pi)^3} E_{\vec{k}} \, \hat{a}_{\vec{k}}^\dagger \hat{a}_{\vec{k}} + \text{const}
\]
where $E_{\vec{k}} = \sqrt{\vec{k}^2 + m^2}$ is the energy of a single mode. (Deriving this from the expression of $T^{00}$ is a standard and highly recommended exercise).

\begin{remark}
    The "$+\text{const}$" term represents the vacuum zero-point energy, which in QFT is technically infinite. We typically eliminate it by defining normal ordering ($:\hat{H}:$) so that the vacuum state has exactly zero energy.
\end{remark}

\subsection{Complex Scalar Field}

The complex scalar field generalizes the real scalar field by introducing a field \(\phi\) that takes complex values. The Lagrangian density for a free complex scalar field is given by:
\[
    \mathcal{L} = |\partial_\mu \phi|^2 - m^2 |\phi|^2 = (\partial_\mu \phi^*)(\partial^\mu \phi) - m^2 \phi^* \phi
\]
A complex field \(\phi\) inherently possesses two independent degrees of freedom. We can parameterize it using two real scalar fields, for instance \(\phi_1\) and \(\phi_2\), such that \(\phi = \frac{1}{\sqrt{2}}(\phi_1 + i\phi_2)\). However, it is mathematically more convenient to treat \(\phi\) and its complex conjugate \(\phi^*\) as the two independent degrees of freedom when applying the Euler-Lagrange equations.

\subsubsection{Continuous Symmetries and Noether's Current}

The Lagrangian \(\mathcal{L}\) is manifestly invariant under a global \(U(1)\) phase transformation. If we rotate the field by a continuous phase \(\alpha\), the transformations are:
\[
    \phi \to e^{i\alpha} \phi, \qquad \phi^* \to e^{-i\alpha} \phi^*
\]
For an infinitesimal transformation (\(\alpha \ll 1\)), the field variations are \(\delta \phi = i\alpha \phi\) and \(\delta \phi^* = -i\alpha \phi^*\). Because this continuous symmetry leaves the action invariant, Noether's theorem dictates the existence of a conserved current \(J^\mu\).

We can derive this current explicitly (solving the suggested exercise) using the standard Noether formula:
\[
    J^\mu = \frac{\partial \mathcal{L}}{\partial(\partial_\mu \phi)} \delta\phi + \frac{\partial \mathcal{L}}{\partial(\partial_\mu \phi^*)} \delta\phi^*
\]
Dropping the overall constant parameter \(\alpha\), and substituting the derivatives \(\frac{\partial \mathcal{L}}{\partial(\partial_\mu \phi)} = \partial^\mu \phi^*\) and \(\frac{\partial \mathcal{L}}{\partial(\partial_\mu \phi^*)} = \partial^\mu \phi\), we obtain:
\[
    J^\mu = (\partial^\mu \phi^*)(i\phi) + (\partial^\mu \phi)(-i\phi^*) = -i (\phi \partial^\mu \phi^* - \phi^* \partial^\mu \phi)
\]
By virtue of the equations of motion (the Klein-Gordon equation for both \(\phi\) and \(\phi^*\)), this current is strictly conserved, meaning \(\partial_\mu J^\mu = 0\).

\subsubsection{Quantization and Antiparticles}

To quantize the theory, we promote the classical fields \(\phi\) and \(\phi^*\) to operator-valued distributions \(\hat{\phi}\) and \(\hat{\phi}^\dagger\). Because the field is complex (\(\hat{\phi} \neq \hat{\phi}^\dagger\)), the Fourier expansion requires two distinct sets of creation and annihilation operators:
\[
    \hat{\phi}(x) = \int \frac{d^3k}{(2\pi)^3} \frac{1}{\sqrt{2E_{\vec{k}}}} \left( \hat{a}_{\vec{k}} e^{-ik \cdot x} + \hat{b}_{\vec{k}}^\dagger e^{ik \cdot x} \right)
\]
\[
    \hat{\phi}^\dagger(x) = \int \frac{d^3k}{(2\pi)^3} \frac{1}{\sqrt{2E_{\vec{k}}}} \left( \hat{a}_{\vec{k}}^\dagger e^{ik \cdot x} + \hat{b}_{\vec{k}} e^{-ik \cdot x} \right)
\]
where \(E_{\vec{k}} = \sqrt{\vec{k}^2 + m^2}\).
The physical interpretation of these operators is fundamental to QFT:
\begin{itemize}
    \item \(\hat{a}_{\vec{k}}\) and \(\hat{a}_{\vec{k}}^\dagger\) annihilate and create \textbf{particles}, respectively.
    \item \(\hat{b}_{\vec{k}}\) and \(\hat{b}_{\vec{k}}^\dagger\) annihilate and create \textbf{anti-particles}, respectively.
\end{itemize}
This implies the existence of two distinct types of particle states. Because they both originate from the same Lagrangian and share the same dispersion relation \(E_{\vec{k}}\), particles and anti-particles possess the exact same mass.

\subsubsection{The Conserved Charge}

The conserved Noether current \(J^\mu\) leads to a conserved Noether charge \(Q = \int d^3x \, J^0\). Upon quantization, promoting the fields to operators and substituting the mode expansions into the integral yields the quantum charge operator \(\hat{Q}\):
\[
    \hat{Q} = \int \frac{d^3k}{(2\pi)^3} \left( \hat{a}_{\vec{k}}^\dagger \hat{a}_{\vec{k}} - \hat{b}_{\vec{k}}^\dagger \hat{b}_{\vec{k}} \right) + \text{const}
\]
The additive constant arises from commutation relations during the derivation and represents the vacuum charge. In practice, we apply \textit{normal ordering} (\(:\hat{Q}:\)) to subtract this infinite vacuum contribution, setting the charge of the vacuum state to zero.

The resulting operator algebraically "counts" the number of particles minus the number of anti-particles. This demonstrates that while particles and anti-particles have the same mass, they carry \textbf{opposite conserved charges} under the \(U(1)\) symmetry.

\subsubsection{The Feynman Propagator}

For the complex scalar field, the causal propagation between two spacetime points is described by the time-ordered expectation value of the field and its conjugate. The Feynman propagator is defined as:
\[
    D_F(x, y) = \langle 0 | T \{ \hat{\phi}(x) \hat{\phi}^\dagger(y) \} | 0 \rangle = \int \frac{d^4k}{(2\pi)^4} \frac{i}{k^2 - m^2 + i0^+} e^{-ik \cdot (x-y)}
\]
Physically, this represents the amplitude for a particle to propagate from \(y\) to \(x\) (if \(x^0 > y^0\)), or for an anti-particle to propagate from \(x\) to \(y\) (if \(y^0 > x^0\)).

\begin{remark}
    The expectation values of two identical fields vanish identically:
    \[
        \langle 0 | T \{ \hat{\phi}(x) \hat{\phi}(y) \} | 0 \rangle = \langle 0 | T \{ \hat{\phi}^\dagger(x) \hat{\phi}^\dagger(y) \} | 0 \rangle = 0
    \]
    This happens because the operator \(\hat{\phi}\hat{\phi}\) changes the net \(U(1)\) charge by two units, meaning it connects the vacuum (charge 0) to a state with charge +2 or -2, making the inner product with the vacuum state strictly zero.
\end{remark}
